\cleardoublepage
\setcounter{section}{0}
\setcounter{table}{0}
\setcounter{figure}{0}
\addcontentsline{toc}{chapter}{Arquiteturas de Integração com Python - \textit{Ryan Faustino Carvalho}}

\begin{center}
    {\fontsize{14}{16}\selectfont\color{sommagreen}\textbf{Arquiteturas de Integração com Python} \\ Estruturação de pipelines de dados} \\[1cm]
    
    {\fontsize{12}{14}\selectfont\textbf{Ryan Faustino Carvalho\textsuperscript{1*}}} \\[0.5cm]
    
    {\fontsize{10}{12}\selectfont
    \textsuperscript{1}Piauí Gov Tech - Teresina, Piauí, Brasil.
    }
\end{center}

\vspace{0.8cm}

\noindent {\fontsize{10}{12}\selectfont\color{sommagreen}\textbf{RESUMO}} \\[0.2cm]
\noindent {\fontsize{10}{12}\selectfont
\textbf{Introdução}: Sistemas escaláveis requerem fundações robustas. \textbf{Objetivo}: Mapear padrões de desenvolvimento de aplicações governamentais. \textbf{Metodologia}: Análise de código e estruturação de microsserviços. \textbf{Resultados e Discussão}: O uso de pipelines automatizados demonstrou maior resiliência. \textbf{Conclusões}: O ecossistema Python provê as ferramentas necessárias para a alta disponibilidade.
}
\vspace{0.5cm}

\begin{tcolorbox}[colframe=sommagreen, colback=white, sharp corners, boxrule=0.8pt, arc=0pt, left=4pt, right=4pt, top=4pt, bottom=4pt]
\fontsize{10}{12}\selectfont
\textcolor{sommagreen}{\textbf{PALAVRAS-CHAVE:}} Django. Next.js. Arquitetura de Sistemas. Backend.
\end{tcolorbox}
\vspace{1cm}

\section{Introdução}
O desenvolvimento de backends robustos requer o entendimento profundo da linguagem escolhida e de suas abstrações mais avançadas \cite{ramalho2022}.

\section{Integração e Deploy}
Em ambientes de segurança pública, a confiabilidade da integração entre ferramentas frontend e a base de dados é fundamental para o sucesso das operações diárias \cite{gomes2025}. A estrutura modular facilita o processo de manutenção e expansão.

\vspace{1cm}
\noindent{\fontsize{12}{14}\bfseries\color{sommagreen}REFERÊNCIAS}
\vspace{0.3cm}

\printbibliography[heading=nossabibliografia]